\documentclass{article}
\usepackage{spconf,amsmath,epsfig,float}
\usepackage{xcolor}

% Example definitions.
% --------------------
\def\x{{\mathbf x}}
\def\L{{\cal L}}

% Title.
% ------
\title{CAN TABKAN MODELS BALANCE ACCURACY AND INTERPRETABILITY IN LIFE INSURANCE RISK ASSESSMENT?}

\name{Tilman Haferbeck, Gian Seifert, Christof Steiner, Cyril Gabriele}
\address{University of St.\ Gallen}

\begin{document}
\maketitle

\begin{abstract}
Life insurance underwriting requires risk models that are both accurate and transparent, since predictions affect pricing and must be justifiable in regulated settings. Traditional actuarial models are interpretable but may miss nonlinear relationships, whereas modern machine learning models often improve predictive performance at the cost of reduced transparency. This paper investigates whether Kolmogorov--Arnold Network (KAN)-based tabular models, in particular TabKAN variants, can improve this accuracy--interpretability trade-off. Using the Prudential Life Insurance Assessment dataset, we compare ChebyKAN, FourierKAN, and BSplineKAN against XGBoost and a multilayer perceptron under a fixed preprocessing and validation protocol. We further analyze sparse KAN configurations to assess whether compact functional structures can retain competitive predictive performance. Our results show that XGBoost remains the strongest dense predictor, while KAN variants consistently outperform the standard neural baseline. In the sparse regime, ChebyKAN provides the most favorable interpretability trade-off by approaching the predictive performance of stronger baselines with a substantially smaller and analytically inspectable structure.
\end{abstract}

\begin{keywords}
TabKAN, Kolmogorov--Arnold Networks, tabular deep learning, life insurance, interpretability
\end{keywords}

\section{Introduction}
\label{sec:intro}

Risk assessment is a central task in life insurance underwriting. Insurers estimate an applicant's mortality or morbidity risk in order to support policy acceptance and pricing decisions, which makes predictive accuracy important for fair premiums and financial stability. At the same time, underwriting is a regulated and high-stakes setting in which model decisions must remain transparent and justifiable to firms, regulators, and policyholders.

Classical actuarial models such as generalized linear models (GLMs) remain attractive because they provide well-understood and interpretable relationships between risk factors and outcomes \cite{Frees2009}. However, their functional forms and interaction effects usually need to be specified manually, which can limit their ability to capture nonlinear patterns in heterogeneous insurance data. Machine learning methods, including tree-based ensembles, can learn such interactions more flexibly and have shown strong results in insurance analytics \cite{Henckaerts2021,Richman2021}. Their drawback is that predictive flexibility often reduces transparency. Post-hoc methods such as SHAP can explain individual predictions, but the underlying model remains complex and only indirectly interpretable \cite{LundbergLee2017,Aas2021}.

This tension is especially relevant for deep learning on tabular insurance data. Although deep neural networks are highly expressive, they are often difficult to interpret and do not consistently outperform strong tabular baselines. Kolmogorov--Arnold Networks (KANs) offer a promising alternative because they replace fixed node activations with learnable univariate functions on network edges \cite{Liu2024}. This design can make learned transformations more explicit while retaining nonlinear approximation power. Recent work has adapted this idea to tabular learning through TabKAN and related architectures \cite{Eslamian2025,Gao2024,Poeta2024}.

However, existing KAN studies focus mainly on generic benchmark datasets. Evidence from domain-specific, high-stakes applications such as life insurance risk assessment is still limited. This paper addresses that gap by asking whether TabKAN models can provide competitive predictive performance while improving model-native interpretability. We make three contributions: we evaluate KAN-based models on a life-insurance underwriting task, compare several TabKAN variants against established tabular baselines, and analyze whether sparse KAN configurations can offer a practical accuracy--interpretability compromise.

\section{Related Work}
\label{sec:related}

Tabular learning is a dominant setting in applied machine learning, particularly in finance, healthcare, and insurance. In this domain, tree-based ensemble methods remain strong baselines because they handle heterogeneous feature types, nonlinear interactions, and medium-sized datasets effectively. Random Forests \cite{Breiman2001}, gradient boosting \cite{Friedman2001}, and specialized implementations such as XGBoost, LightGBM, and CatBoost \cite{ChenGuestrin2016,Ke2017,Prokhorenkova2018} have therefore become standard references for structured prediction. By contrast, conventional multilayer perceptrons (MLPs) often fail to consistently outperform tree-based methods on tabular data \cite{ShwartzZiv2022,Grinsztajn2022}. Recent neural tabular models such as TabNet, TabTransformer, FT-Transformer, NODE, and DeepGBM introduce attention, differentiable trees, or hybrid inductive biases \cite{ArikPfister2021,Huang2020,Gorishniy2021,Popov2020,Ke2019}. Nevertheless, broad empirical comparisons still show that new neural architectures should be evaluated against strong boosting baselines rather than only against MLPs \cite{ShwartzZiv2022,Grinsztajn2022}.

Insurance analytics adds a second requirement: interpretability. Predictive models are used for underwriting, pricing, claims prediction, and fraud detection, where decisions can have financial and regulatory consequences. GLMs remain common in actuarial work because their coefficients are directly interpretable \cite{Frees2009}. More flexible machine learning models can improve predictive accuracy in insurance applications \cite{Henckaerts2021,Richman2021}, but their complexity raises questions about accountability and bias. Explainable Artificial Intelligence (XAI) methods, especially SHAP, are therefore widely used to attribute predictions to input features \cite{LundbergLee2017,Guidotti2018}. Such explanations are useful, but they remain post-hoc: they approximate or summarize the behavior of a model whose internal decision structure is not itself transparent.

KANs have recently been proposed as a neural architecture inspired by the Kolmogorov--Arnold representation theorem \cite{Kolmogorov1957,Liu2024}. Instead of applying fixed nonlinearities at neuron outputs, KANs learn univariate functions on edges, which can expose the functional transformations learned by the model. TabKAN extends this idea to tabular data and reports competitive classification performance \cite{Eslamian2025}; related work combines KAN-style components with transformer-based tabular architectures \cite{Gao2024}. These properties make KANs relevant for insurance, where data are predominantly tabular and transparency is valuable. Yet their suitability for life-insurance risk assessment remains underexplored, motivating our empirical comparison of TabKAN variants with established baselines. Beyond architectural innovations, recent work introduces tabular foundation models that are pretrained on synthetic priors and applied to new tasks via in-context learning rather than per-task gradient updates; TabPFN \textcolor{red}{[tabpfn reference to add]} is the principal example and serves here as a contemporary no-tuning baseline.

\section{Approach and Experimental Setup}
\label{sec:method}

\begin{figure}[H]
    \centering
    \includegraphics[width=\columnwidth]{images/methodology.png}
    \caption{Experimental workflow.}
    \label{fig:methodology}
\end{figure}

\subsection{Dataset and Prediction Task}

We evaluate all models on the \textit{Prudential Life Insurance Assessment} dataset, a public benchmark for insurance risk prediction. The dataset contains 59{,}381 training instances and 19{,}765 test instances with 128 tabular features covering demographic information, biometric measurements, product attributes, and medical history variables. The target variable, \textit{Response}, is an ordinal label from 1 to 8, where higher values indicate higher underwriting risk. We therefore treat the task as ordinal prediction rather than standard multiclass classification.

All experiments use the preprocessing pipeline of the reproduced baseline paper. Keeping preprocessing fixed allows us to focus on model architecture rather than feature-engineering differences. Only model-required technical adjustments are added, most importantly explicit NaN handling for neural models that cannot process missing values directly. All models are trained and validated on identical 80/20 splits.

\subsection{Models and Training Protocol}

We compare one strong tree-based baseline, one conventional neural baseline, and three KAN-based models. XGBoost serves as the main tabular benchmark, while a standard MLP provides a neural reference. The KAN variants differ in the basis functions used for their learned edge transformations: ChebyKAN uses Chebyshev polynomial expansions, FourierKAN uses Fourier basis functions, and BSplineKAN uses B-spline basis functions. This selection allows us to test whether different KAN parameterizations lead to different accuracy--interpretability trade-offs. We additionally include TabPFN, a transformer-based tabular foundation model, as a no-tuning reference baseline. TabPFN is run in regression mode on a deterministic 10{,}000-row training subsample (per seed) using the published in-context inference path, and the same QWK-based threshold calibration is applied as for the MLP and KAN models.

For the MLP and all KAN variants, we implement the task as ordinal regression. Each model is trained with an MSE regression head that outputs a continuous score; validation-based thresholds then map scores to the eight discrete response classes. Thresholds are tuned to maximize Quadratic Weighted Kappa (QWK), which is the primary evaluation metric and penalizes larger ordinal errors more strongly than smaller ones. XGBoost is trained in the corresponding baseline regression setup with the same QWK-based threshold calibration.

\subsection{Hyperparameter and Sparsity Optimization}

Hyperparameter optimization is performed with Optuna. In the first stage, each model is optimized for validation QWK using 100--151 trials. The search spaces include standard optimization and capacity parameters such as learning rate, depth, hidden width, batch size, dropout, weight decay, and, for XGBoost, tree depth, subsampling, number of estimators, and regularization terms. For KANs, we additionally tune architecture-specific basis parameters such as polynomial degree, Fourier grid size, and spline-grid settings.

In the second stage, we conduct multi-objective Optuna studies for the KAN models to jointly analyze validation QWK and sparsity. Sparsity is used as a proxy for structural interpretability: a model with fewer active edges is easier to inspect and, for certain KAN parameterizations, easier to express analytically. This two-stage protocol first identifies strong predictive configurations and then evaluates whether KANs can retain performance under increasingly compact model structures.

\section{Results}
\label{sec:results}

\subsection{Predictive Performance}
\label{sec:results_optuna}

Table~\ref{tab:single_objective_optuna} reports the best validation QWK from the single-objective Optuna studies. XGBoost achieves the strongest overall result with a QWK of 0.6546. Among the KAN variants, FourierKAN performs best (0.6410), followed by BSplineKAN (0.6320) and ChebyKAN (0.6254). All KAN variants outperform the MLP baseline (0.6129), indicating that KAN parameterizations provide a useful neural inductive bias for this tabular task, although they do not surpass the tree-based baseline in the dense predictive setting.

\begin{table}[H]
\centering
\caption{Best validation QWK from single-objective Optuna tuning. For TabPFN, no Optuna study is performed; the reported value is the mean validation QWK across three seeds at default hyperparameters.}
\label{tab:single_objective_optuna}
\begin{tabular}{lcc}
\hline
Model & Trials & Best QWK \\
\hline
XGBoost   & 100 & 0.6546 \\
TabPFN    & --  & 0.6442 \\
MLP       & 100 & 0.6129 \\
ChebyKAN  & 151 & 0.6254 \\
FourierKAN& 150 & 0.6410 \\
BSplineKAN& 150 & 0.6320 \\
\hline
\end{tabular}
\end{table}

The multi-objective studies show that predictive performance and sparsity trade off differently across KAN variants. ChebyKAN reaches a Pareto-optimal QWK of 0.6061 at sparsity 0.8534 and still retains a QWK of 0.5124 near maximal sparsity (0.9988). FourierKAN has a steeper trade-off: its best Pareto point reaches 0.5875 at sparsity 0.5974, but near-maximal sparsity reduces QWK to 0.1465. Thus, while FourierKAN is the strongest KAN in the dense setting, ChebyKAN provides the more favorable sparsity--accuracy frontier.

\subsection{Interpretability and Model Analysis}
\label{sec:interpretability}

We complement the validation results with held-out outer-test QWK for selected dense and sparse configurations (Table~\ref{tab:interpretability_outer_test}). The full-feature rows compare dense models, while the 20-feature rows evaluate matched sparse settings. For the latter, each model uses its own top-20 features, selected by native coefficient importance for KANs and by mean-$|\mathrm{SHAP}|$ for XGBoost. Each 20-feature configuration is then tuned with a 50-trial Optuna search. For KANs, interpretability depends not only on per-edge symbolic recovery from basis coefficients, but also on how many edges remain active and whether the composed model can be written as a compact expression.

\begin{table}[H]
  \centering
  \caption{Held-out outer-test QWK and structural size. Full-feature intervals are 95\% t-intervals across three seeds; sparse-regime intervals are 95\% bootstrap confidence intervals
  over the outer-test set.}
  \label{tab:interpretability_outer_test}
  \scriptsize
  \begin{tabular}{lccl}
  \hline
  Model & QWK & Size & Explanation \\
  \hline
  \multicolumn{4}{l}{\emph{Full-feature regime}} \\
  XGBoost & 0.648 [0.628, 0.669] & 950 trees & SHAP \\
  TabPFN, default & 0.636 [0.629, 0.644] & in-context & permutation \\
  ChebyKAN, pruned Pareto & 0.624 [0.615, 0.633] & 17{,}482 active edges & native edges \\
  FourierKAN, pruned Pareto & 0.630 [0.602, 0.658] & 17{,}855 active edges & native edges \\
  \hline
  \multicolumn{4}{l}{\emph{Sparse regime (20 features)}} \\
  XGBoost, tuned & 0.613 [0.599, 0.625] & 950 trees & SHAP \\
  TabPFN, default & 0.562 [0.547, 0.575] & in-context & permutation \\
  ChebyKAN, sparse & 0.593 [0.579, 0.607] & 2{,}761 edges & closed-form model\\
  FourierKAN, sparse & 0.577 [0.564, 0.591] & 467 edges & closed-form model \\
  \hline
  \end{tabular}
\end{table}

\subsubsection{Full-feature regime: KANs vs. XGBoost + SHAP}
\label{sec:interpretability_dense}
% CYRIL INTERPRETABILITY START 

In the full-feature interpretability setting, the models differ not only in predictive structure but also in the type of explanation they provide. For XGBoost, interpretability is obtained post hoc using predicted-class Tree SHAP values. Global feature importance is computed from mean-$|\mathrm{SHAP}|$, while feature-effect plots show applicant-level SHAP contributions for the class predicted by the model. For the KAN variants, we use model-native diagnostics from the tuned ChebyKAN and FourierKAN modules: feature rankings are derived from native coefficient/edge importance, and feature effects are summarized by partial-dependence plots (PDPs) computed directly from the KAN modules. Figure~\ref{fig:feature_effect_comparison} compares these KAN PDPs with the corresponding XGBoost SHAP dependence plots.

The three explanation methods agree on a meaningful core of underwriting variables, but they do not induce identical rankings. BMI is the highest-ranked feature for XGBoost and ChebyKAN, and it is ranked second by FourierKAN. Seven variables appear in the top 20 for all three methods. Relative to XGBoost SHAP, ChebyKAN shares twelve top-20 features with a rank-scaled overlap of 5.1675, whereas FourierKAN shares nine top-20 features with a rank-scaled overlap of 4.1725. Thus, the models recover a shared signal core, while the remaining rank differences indicate architecture-specific representations rather than a single universal explanation.

The feature-effect panels in Figure~\ref{fig:feature_effect_comparison} cover four representative features that are important across the models and include both continuous and binary inputs: BMI, Medical\_History\_15, Product\_Info\_4, and Medical\_Keyword\_3. The XGBoost row shows the distribution of local SHAP contributions across applicants, with a binned or state-wise mean overlay. The KAN rows show global PDP responses, i.e., the average predicted risk score obtained when one feature is swept while all other applicant features are held at their observed values. These vertical axes should not be compared numerically: SHAP values are local additive contributions to a class-specific prediction, whereas KAN PDP curves are global marginal summaries of the predicted risk score.

\begin{figure}[t]
    \centering
    \includegraphics[width=0.49\textwidth]{images/feature_effect_comparison.png}
    \caption{
    Feature-level explanation comparison between the XGBoost baseline and the tuned KAN models. Columns show four features selected for high cross-model relevance and mixed feature semantics. The XGBoost row reports predicted-class SHAP values for individual applicants with a binned or state-wise mean overlay. The ChebyKAN and FourierKAN rows report model-native PDP curves. The panels compare explanation geometry rather than a common vertical scale: SHAP values are local additive contributions, whereas KAN PDPs are global marginal responses of the predicted risk score.
    }
    \label{fig:feature_effect_comparison}
\end{figure}

The panels highlight the complementary strengths of the two interpretability approaches. XGBoost + SHAP provides mature local explanations: each point corresponds to one applicant and can be interpreted as a class-specific contribution to that applicant's prediction. This is useful for case-level underwriting review, but the explanation remains post hoc because it summarizes the behavior of a trained tree ensemble through an additional attribution framework. The KAN PDPs, by contrast, expose learned response functions directly from the model structure. For BMI, XGBoost and ChebyKAN show a broadly decreasing response at higher BMI values, whereas FourierKAN displays a more oscillatory nonlinear response. For Medical\_History\_15, the models agree on feature relevance but differ in effect geometry: XGBoost shows an increasing SHAP trend, while the KAN PDPs show a sharp low-end rise followed by a decline across much of the range. Product\_Info\_4 provides one of the strongest examples of cross-model rank agreement, with all three models ranking it fifth and all curves showing nonlinear behavior. Medical\_Keyword\_3 illustrates the discrete binary case: XGBoost and ChebyKAN show a strong change between the observed states, whereas FourierKAN assigns high importance but a smaller marginal change.

In predictive terms, XGBoost remains the strongest baseline in this dense setting, reaching a QWK of 0.6455. FourierKAN is close, with a QWK of 0.6331 before pruning and 0.6332 for the interpreted module, while ChebyKAN reaches a QWK of 0.6251 before pruning and 0.6161 for the interpreted module. The KAN variants therefore operate in the same broad predictive regime as XGBoost, especially FourierKAN, but they do not outperform the XGBoost baseline on absolute QWK.

TabPFN reaches an outer-test QWK of $0.636$ in the full-feature setting, comparable to but not surpassing XGBoost. Its top-5 features (by mean-replacement ablation impact on QWK) overlap by $3$ with XGBoost-SHAP and by $2$ with the ChebyKAN native ranking, indicating that the foundation model recovers the same core underwriting signal. TabPFN interpretability is similar in kind to XGBoost: both rely on post-hoc methods (SHAP, ablation, permutation importance) and neither admits a model-native closed-form representation; TabPFN supports attention-based attribution over training applicants in principle, but extracting it from the in-context inference path requires hooks beyond the standard scikit-learn interface and is left to future work, so it remains on the post-hoc side of our comparison.

The interpretability advantage of the tuned KANs should be framed carefully. The tuned ChebyKAN and FourierKAN models are dense rather than compact symbolic models: ChebyKAN retains 17{,}482 active edges, and FourierKAN retains 17{,}855 active edges. Both KAN variants achieve mean per-edge symbolic recovery $R^2=1.0$, indicating that individual learned edge functions can be faithfully represented in their native basis. However, this is an edge-level symbolic fidelity result, not a claim that the full deployed network is a simple end-to-end closed-form expression. The composition of many active edges and architectural components makes the full model substantially more complex than any individual edge function.

Overall, XGBoost + SHAP and KAN-native explanations provide different but complementary forms of transparency. SHAP is better suited for local additive explanations of a strong black-box tree ensemble. KAN explanations are more model-native and make learned feature-response geometry directly inspectable without requiring a separate attribution framework. In this full-feature dense comparison, the main advantage of the KAN variants is therefore not superior accuracy over XGBoost, but the availability of native feature-effect diagnostics from models that remain competitive in predictive performance.
% CYRIL INTERPRETABILITY END

\subsubsection{Sparse regime: compact ChebyKAN explanations}
\label{sec:interpretability_sparse}

The 20-feature results show a different pattern. The 95\% bootstrap confidence intervals form a partial ordering: XGBoost (0.613) overlaps the sparse ChebyKAN (0.593), the sparse ChebyKAN overlaps the sparse FourierKAN (0.577), but XGBoost and FourierKAN do not overlap. Both KAN configurations are obtained in two steps: a 50-trial Optuna search at the 20-feature budget tunes architecture and optimiser with $\lambda=0$, then a 33-trial $\lambda$ Pareto sweep on the winner architecture picks a strong sparse-reference $\lambda$ at sparsity $\geq 0.85$. The sparse ChebyKAN has two hidden layers (depth~2, width~128), degree 7, $\lambda=0.00236$; the sparse FourierKAN has one hidden layer (depth~1, width~128), grid size 8, $\lambda=0.0005$. Both turn LayerNorm off and prune edges whose mean activation magnitude falls below $2\times10^{-3}$. This leaves 2{,}761 of 18{,}944 ChebyKAN edges (approximately 85\% pruning) and 467 of 2{,}560 FourierKAN edges (approximately 82\% pruning), while preserving outer-test performance within the bootstrap uncertainty of the stronger baselines. In the same 20-feature regime, TabPFN on its own ablation-derived top-20 reaches outer-test QWK $0.562$, $-0.051$ from XGBoost on its top-20, indicating that the foundation model relies more heavily on the full feature set than the sparse KAN variants do.

This compact structure is central to the interpretability claim. Each active ChebyKAN edge is a polynomial of degree at most 7 in $\tanh(x)$, read directly from the learned Chebyshev coefficients; for example, one learned BMI edge is $-1.926\,T_7(\tanh x)+1.134\,T_4(\tanh x)+1.122\,T_5(\tanh x)$. The sparse FourierKAN admits a parallel reading: each active edge is a finite Fourier series in $\tanh(x)$ with up to eight sine and eight cosine harmonics, read directly from the learned Fourier coefficients (R$^2 = 1.000$ basis-native, as for ChebyKAN). Because LayerNorm is removed and the FourierKAN winner is only one hidden layer deep, the active edges of either model compose through the KAN layers into a single closed-form expression over the retained inputs. The FourierKAN closed-form property is, however, conditional on this depth-1 configuration: composition of finite Fourier series through nonlinear edges does not in general remain a finite Fourier series, so deeper FourierKANs would only admit a nested closed form. ChebyKAN preserves this property at any depth, since a polynomial of polynomials is again a polynomial; for this reason ChebyKAN is used here as the more general illustration of model-native interpretability. The same expression supports exact symbolic derivatives, with the ChebyKAN BMI derivative for applicant 55728 matching PyTorch autograd up to numerical precision. Unlike SHAP explanations for XGBoost, these explanations are therefore model-native rather than post-hoc.

\begin{figure}[H]
    \centering
    \includegraphics[width=\columnwidth]{images/fig1_interpretability.pdf}
    \caption{Analytic per-feature decomposition of the sparse ChebyKAN prediction for applicant 55728. Integrated-gradient contributions along the closed-form polynomial path sum to the difference between the reference score and the applicant score.}
    \label{fig:fig1_interpretability}
\end{figure}

Figure~\ref{fig:fig1_interpretability} illustrates the resulting local explanation. Starting from a low-risk reference applicant with score $-2.10$, the integrated-gradient contributions of the retained features sum to the applicant's predicted score of 3.57, corresponding to class 3. The largest positive contributions come from medical-history and medical-keyword variables, while age and product-related features lower the score. This example shows how the sparse ChebyKAN can provide both a prediction and an analytically grounded decomposition of the same prediction. By contrast, TabPFN can score the same applicant but, in our setup, exposes neither a function-level decomposition nor a directly readable case-based justification through standard interfaces, so the closed-form path of Figure~\ref{fig:fig1_interpretability} has no equivalent on the foundation-model side of the comparison.

Overall, TabKAN occupies a distinct point on the accuracy--interpretability frontier. XGBoost remains the strongest dense baseline, and FourierKAN is the strongest dense KAN. In the sparse 20-feature regime, both KAN variants admit model-native closed-form expressions: ChebyKAN reaches higher QWK at 2{,}761 active edges, while FourierKAN reaches a more compact 467-edge structure at slightly lower QWK. Both remain close to the best 20-feature models in predictive performance.

\section{Discussion and Conclusion}
\label{sec:discussion}

This study asked whether TabKAN models can balance accuracy and interpretability in life insurance risk assessment. The answer is nuanced. In terms of pure predictive performance, XGBoost remains the strongest model, confirming the continued strength of tree-based ensembles on tabular data. However, all KAN variants outperform the standard MLP, which suggests that KAN-based parameterizations are more suitable than conventional neural networks for this underwriting task. Among the dense KANs, FourierKAN achieves the highest QWK, but its advantage weakens when sparsity is imposed.

The main interpretability result is therefore not that TabKAN replaces XGBoost as the most accurate model, but that both sparse KAN variants reach competitive regions of the accuracy--interpretability frontier with closed-form structure. The sparse ChebyKAN (2{,}761 active edges) overlaps the confidence interval of the strongest 20-feature baseline; the sparse FourierKAN reaches the same closed-form properties at only 467 active edges, at a slightly lower QWK. Both allow closed-form edge recovery, analytic feature decompositions, and exact symbolic derivatives. This provides a stronger form of model-native transparency than post-hoc explanations alone and is particularly relevant in regulated insurance settings, where compact and inspectable models are valuable. Even when extending the comparison to a contemporary tabular foundation model, the relative position of the sparse KAN configurations is preserved: TabPFN provides accuracy in the same broad regime as XGBoost without offering model-native interpretability, leaving the sparse ChebyKAN as the only configuration in the comparison that combines competitive 20-feature accuracy with closed-form structure.

Several limitations remain. The empirical evidence is based on a single public insurance dataset, and the fixed preprocessing pipeline may influence the relative model ranking. Sparsity is also only a proxy for interpretability; future work should evaluate whether actuaries or regulators actually find the resulting KAN explanations useful. Finally, KAN training and hyperparameter optimization are computationally more involved than standard XGBoost workflows. Future research should therefore test TabKAN models on additional insurance datasets, explore ordinal losses tailored directly to QWK, and compare model-native KAN explanations with post-hoc XAI methods in user-centered evaluation settings. TabPFN was used at default hyperparameters and on a deterministic 10{,}000-row training subsample as a no-tuning reference; per-task fine-tuning or larger context configurations could change the predictive ranking but lie outside the scope of the present comparison.

\bibliographystyle{IEEEbib}
\bibliography{references}

\end{document}
